\section{Additive distribution of the conductor weights}
\label{sec:additive}

The factors \(f_p^*\) introduced in Section~\ref{sec:exterior} retain the
local lattice condition of an order at the primes in \(S\).  We now average their
product on the real trace plane.  The estimate will also permit a residue
class at a small modulus prime to \(S\); this is the form needed to average
the factors at the remaining primes in Section~\ref{sec:sieve}.

Fix a nonzero integer \(m\).  In the three real embeddings of \(F\), write
\(z=(z_1,z_2,z_3)\), so that \(\Tr z=z_1+z_2+z_3\).  Choose two distinct
indices \(i,j\), a fixed nonnegative function
\(\Phi\in C_c^\infty(\R^2)\) which is at least one on \([-1,1]^2\), and a
fixed nonnegative \(\chi\in C_c^\infty(\R)\) with
\[
 \chi(m)=1,\qquad
 \supp\chi\subset(m-\tfrac13,m+\tfrac13).
\]
For \(c=(c_i,c_j)\in\R^2\) and positive \(r_i,r_j\), put
\begin{equation}
 \phi(z)=
 \Phi\left(\frac{z_i-c_i}{r_i},\frac{z_j-c_j}{r_j}\right)\chi(\Tr z),
 \qquad
 A_\phi=r_i r_j\int_{\R^2}\Phi(u,v)\,du\,dv .
 \label{eq:add-box}
\end{equation}
Thus \(A_\phi\) is the integral of \(\phi\) on \(\Tr z=m\), using \(z_i,z_j\)
as coordinates on that plane.  Fix \(r_*>0\), keep \(r_i,r_j\le r_*\), and
require the supports of their planar factors to lie in a fixed
bounded region.  The supports of all the resulting \(\phi\)'s then lie in a
fixed compact subset \(\Omega_\infty\) of \(\R^3\) on which \(\Tr z\) stays
away from zero.  All constants below may depend on \(m,\Phi,\chi\), the
upper bound for the scales, and this bounded region.  They do not depend on
the locations \(c\) or on the scales within their specified range.

For \(p\notin S\), the algebra \(F\otimes\Q_p\) is unramified and its
integral trace map \(R_{F,p}\to\Z_p\) is onto.  The affine set
\[
 \mathcal T_{m,p}=\{z\in R_{F,p}:\Tr z=m\}
\]
is consequently a translate of a free rank-two direct summand of \(R_{F,p}\).
A residue class on this set modulo \(p^a\) means a translate of
\(p^a\ker(\Tr:R_{F,p}\to\Z_p)\).  Equivalently, if \(z_p^0\in\mathcal T_{m,p}\)
represents the class, its members satisfy \(z-z_p^0\in p^aR_{F,p}\).
There are \(p^{2a}\) such classes.

\begin{proposition}[Additive distribution]
\label{prop:additive-equidistribution}
There are absolute constants \(\tau,\theta>0\) with the following property.
Let \(K,\OO,F\) be as in Section~\ref{sec:exterior}, put
\(X=[\OO_K:\OO]\sqrt{|\Disc(K)|}\), and take \(S\) as in that section for
the fixed nonzero integer \(m\).  In Equation~\eqref{eq:add-box}, choose any
two real partition coordinates and scales
\[
 X^{-\tau}\le r_i,r_j\le r_*
\]
whose planar supports stay in the fixed bounded region above.  Let
\(1\le l\le X^\theta\) be prime to every prime in \(S\).  For each \(p\mid l\)
choose any \(z_p^0\in\mathcal T_{m,p}\), specifying a class modulo
\(p^{v_p(l)}\).  Then
\begin{equation}
 \sum_{\substack{z\in F,\ \Tr z=m\\
                  z\in R_{F,p}\ (p\notin S)\\
                  z-z_p^0\in p^{v_p(l)}R_{F,p}\ (p\mid l)}}
       \phi(z)\prod_{p\in S}f_p^*(z)
 =
 A_\phi\,l^{-2}\prod_{p\in S}d_p(m)
       \bigl(1+O(X^{-\theta})\bigr).
 \label{eq:add-equidistribution}
\end{equation}
Here \(d_p\) is the trace density of Lemma~\ref{lem:fourier-scalar-density}.
The error is uniform over all the fields, orders, ordered real embeddings,
permitted box locations and scales, moduli \(l\), and chosen residue classes.
The implied constant and the lower threshold for \(X\) depend only on the
fixed real data listed above and on \(m\).
\end{proposition}

\subsection{Finite expansion and Poisson summation}

We first assemble the local inverse transforms into a finite global
expansion.  This is the point at which the fact that the cubic algebra is a
field is essential.  Let \(\mathfrak d_F\) be the different of \(F\), and
write \(\mathfrak d_p=\mathfrak d_F R_{F,p}\) for its localization.
For a choice of denominator shells at the primes of \(S\), write
\[
 \mathfrak h=\prod_{p\in S}\mathfrak h_p,\qquad
 H=\prod_{p\in S}H_p=\Nm\mathfrak h,\qquad
 G_0=\prod_{p\in S}p^{g_p}.
\]
Local ideals here and below are identified with global ideals supported at
the corresponding primes.  The functions \(K_{p,\mathfrak h_p}\) and their
boundary sets \(\partial_{\mathfrak h_p}(z)\) are those of
Lemma~\ref{lem:fourier-shell-estimates}.  Set
\[
 \mathscr L_{\mathfrak h}=\{p\in S:g_p>0\text{ or }p\le5\}.
\]

\begin{lemma}[Global bounds for denominator profiles]
\label{lem:add-profile-bounds}
For every profile one has
\begin{equation}
 G_0\le [R_F:C_F]\le H G_0.
 \label{eq:add-profile-indices}
\end{equation}
Suppose that \(z\in F\) is integral outside \(S\), its real image lies in
\(\Omega_\infty\), and
\(\prod_{p\in S}K_{p,\mathfrak h_p}(z)\ne0\).  Define
\[
 T_\partial=T_\partial(\mathfrak h,z)
   =\prod_{\substack{p\in S\setminus\mathscr L_{\mathfrak h}\\
                     \mathfrak q\in\partial_{\mathfrak h_p}(z)}}
       \Nm\mathfrak q ,
\]
where an empty product is one.  There is a constant depending only on
\(\Omega_\infty\) such that
\begin{align}
 H&\ll D_FG_0^{24}T_\partial,
 &T_\partial&\le\left(\prod_{p\in S}p\right)^3,
 \label{eq:add-profile-norm}\\
 \left|\prod_{p\in S}K_{p,\mathfrak h_p}(z)\right|
   &\ll 8^{|S|}T_\partial^{-1/2}.
 \label{eq:add-profile-inverse}
\end{align}
All such \(z\) lie in a fractional \(R_F\)-ideal
\(\mathfrak a_{\mathfrak h}\), independent of \(z\), with
\begin{equation}
 \Nm(\mathfrak a_{\mathfrak h})^{-1}
       \le 30^{12}D_FG_0^{24}.
 \label{eq:add-profile-support}
\end{equation}

Set \(C_3=371\).  For all sufficiently large \(X\), depending only on the
fixed real data and \(m\), a profile contributing at any such \(z\) has
\(H\le X^{C_3}\).  For these \(X\) and every \(z\in F\) integral outside
\(S\) with real image in \(\Omega_\infty\),
\begin{equation}
 \prod_{p\in S}f_p^*(z)
  =\sum_{\substack{\mathfrak h=(\mathfrak h_p)_{p\in S}\\H\le X^{C_3}}}
        \prod_{p\in S}K_{p,\mathfrak h_p}(z).
 \label{eq:add-profile-expansion}
\end{equation}
The number of profiles in this sum is \(O_\eta(X^\eta)\) for every fixed
\(\eta>0\), uniformly in the arithmetic data.
\end{lemma}

\begin{proof}
Let \(M_p=\mathfrak h_p\cap C_{F,p}\).  The additive lattice identity
\[
 [R_{F,p}:C_{F,p}]
   =[R_{F,p}:C_{F,p}+\mathfrak h_p]\,[\mathfrak h_p:M_p]
\]
has first factor between \(1\) and
\([R_{F,p}:\mathfrak h_p]=H_p\).  The index of \(C_F\) is supported on \(S\)
by Lemma~\ref{lem:ext-order-scale}; multiplying gives
Equation~\eqref{eq:add-profile-indices}.

At \(p\in\mathscr L_{\mathfrak h}\), the local support is contained in
\(p^{-4(g_p+1)}\mathfrak h_p\mathfrak d_p^{-1}\).
At every other \(p\in S\), the support loses at most one power of each
boundary prime \(\mathfrak q\), and loses no power in the other components.
It follows that a supporting \(z\) belongs to
\[
 \mathfrak b_{\mathfrak h,z}
  =\mathfrak h\mathfrak d_F^{-1}
     \prod_{p\in\mathscr L_{\mathfrak h}}
           (pR_F)^{-4(g_p+1)}
     \prod_{\substack{p\in S\setminus\mathscr L_{\mathfrak h}\\
                      \mathfrak q\in\partial_{\mathfrak h_p}(z)}}
           \mathfrak q^{-1}.
\]
The real trace condition excludes \(z=0\).  Since \(F\) is a field, its norm
is nonzero, and containment of the principal ideal \((z)\) in
\(\mathfrak b_{\mathfrak h,z}\) gives
\[
 |\Nm_{F/\Q}z|
   =\Nm(\mathfrak b_{\mathfrak h,z})
        [\mathfrak b_{\mathfrak h,z}:(z)]
   \ge \frac{H}
      {D_FT_\partial\prod_{p\in\mathscr L_{\mathfrak h}}p^{12(g_p+1)}}.
\]
The left side is bounded above on \(\Omega_\infty\).  For \(g_p>0\) we have
\(g_p+1\le2g_p\); the only additional factors, for \(g_p=0\), are at
\(2,3,5\).  Hence
\[
 \prod_{p\in\mathscr L_{\mathfrak h}}p^{12(g_p+1)}
       \le30^{12}G_0^{24}.
\]
This proves the first inequality of Equation~\eqref{eq:add-profile-norm}.
Each boundary prime occurs once, and
\(\sum_{\mathfrak q\mid p}f_{\mathfrak q}\le
\sum_{\mathfrak q\mid p}e_{\mathfrak q}f_{\mathfrak q}=3\);
this proves the bound for \(T_\partial\).  Equation~\eqref{eq:add-profile-inverse}
is the product of the boundary bounds of
Lemma~\ref{lem:fourier-shell-estimates}; its absolute implied constant
accounts for the three small primes.

Dropping \(\mathfrak h_p\) in the first support condition enlarges that
support.  In the second, the exponent \(v_{\mathfrak q}(\mathfrak h_p)-1\)
at a boundary is nonnegative.  Thus every supporting
point belongs to the ideal
\[
 \mathfrak a_{\mathfrak h}
   =\mathfrak d_F^{-1}
      \prod_{p\in\mathscr L_{\mathfrak h}}
             (pR_F)^{-4(g_p+1)} ,
\]
which gives Equation~\eqref{eq:add-profile-support}.

Lemma~\ref{lem:ext-order-scale} gives
\(\Delta_F\ll X^{28}\),
\(D_F\le\Delta_F\), and
\([R_F:C_F]\le\Delta_F^{1/2}\).  The definition of \(S\) gives
\(\prod_{p\in S}p\ll_m X^2\).  Equations
\(\eqref{eq:add-profile-indices}\)--\(\eqref{eq:add-profile-norm}\) therefore
give \(H\ll X^{28+24\cdot14+6}=X^{370}\).  The fixed implied constant is
absorbed by \(X^{371}\) for sufficiently large \(X\).

A nonzero field element has nonzero image in every component of
\(F\otimes\Q_p\).  Lemma~\ref{lem:fourier-inverse-expansion} therefore
identifies the density \(f_p^*(z)\) with its pointwise finite sum of inverse
shells at all the present \(z\).  The norm bound just proved excludes every
product with \(H>X^{371}\), uniformly in \(z\), and proves
Equation~\eqref{eq:add-profile-expansion}.  Finally, for any fixed \(u>0\),
Rankin's inequality bounds the number of profiles with \(H\le X^{C_3}\) by
\[
 X^{C_3u}\prod_{p\in S}(1-p^{-u})^{-3}.
\]
There are at most three prime-ideal components over each \(p\).  The radical
bound on \(S\) implies \(|S|=O(\log X/\log\log X)\), since the product of
\(k\) distinct primes is at least \(k!\).  The displayed Euler product is
then \(X^{o(1)}\) for fixed \(u\).  Taking \(u\) arbitrarily small proves
the asserted bound for every \(\eta>0\).
\end{proof}

We apply this expansion before imposing Poisson summation.  For \(p\notin S\)
let
\[
 \eta_{p,l}(z)=
 \begin{cases}
  \mathbf 1_{z_p^0+p^{v_p(l)}R_{F,p}}(z),&p\mid l,\\
  \mathbf 1_{R_{F,p}}(z),&p\nmid l .
 \end{cases}
\]
Let \(\eta_{\mathfrak h,l}\) be the finite-adelic test function whose local
factor is \(K_{p,\mathfrak h_p}\) at \(p\in S\) and \(\eta_{p,l}\) elsewhere.
By Lemma~\ref{lem:fourier-shell-estimates}, every \(K_{p,\mathfrak h_p}\)
is supported on elements of integral trace.  The factors outside \(S\) are
supported in \(R_{F,p}\) and also have integral trace.  Thus any \(z\in F\)
with \(\eta_{\mathfrak h,l}(z)\ne0\) has rational trace integral at every
finite prime, hence an integer.  Only \(m\) meets \(\supp\chi\).  The finite
profile expansion in Lemma~\ref{lem:add-profile-bounds} consequently shows
that the left side of
Equation~\eqref{eq:add-equidistribution}, denoted by \(\Sigma\), equals
\begin{equation}
 \Sigma
   =\sum_{\substack{\mathfrak h\\H\le X^{C_3}}}
       S_{\mathfrak h},\qquad
 S_{\mathfrak h}
   =\sum_{z\in F}\phi(z)\eta_{\mathfrak h,l}(z).
 \label{eq:add-original-profiles}
\end{equation}
Here \(C_3=371\) is fixed.  The preceding lemma makes the profile interchange
finite.  Each inner sum also has finitely many terms, since its finite
conditions place \(z\) in a fractional ideal and its real support is compact.

In the local results of Section~\ref{sec:local-fourier}, choose the finite
components of the standard character of \(\mathbb A_\Q/\Q\); they all have
conductor \(\Z_p\).  Compose this global character with the field trace.
At a finite place the self-dual measure satisfies
\(\vol(R_{F,p})=|D_F|_p^{1/2}\), while at infinity it is Lebesgue measure
in the three embeddings.  The real covolume of \(R_F\) is \(\sqrt{D_F}\)
and \(\prod_p\vol(R_{F,p})=D_F^{-1/2}\).  Thus the diagonal copy of \(F\)
has covolume one for these measures.  Write
\[
 \widehat\phi(A)=\int_{\R^3}\phi(z)\psi_\infty(\Tr(Az))\,dz
\]
for the real transform.  Field Poisson summation in this
normalization \cite[Section~4.2, Lemma~4.2.4]{Tate1967} gives
\begin{equation}
 S_{\mathfrak h}
   =\sum_{A\in F}\widehat\phi(A)
             \widehat\eta_{\mathfrak h,l}(A),\qquad
 \widehat\eta_{\mathfrak h,l}(A)
   =\prod_{p\in S}\gamma_{p,\mathfrak h_p}(A)
       \prod_{p\notin S}\widehat\eta_{p,l}(A).
 \label{eq:add-poisson}
\end{equation}
Here \(\gamma_{p,\mathfrak h_p}\) is the shell restriction of \(\gamma_p\),
extended by zero, and hats use the chosen character.  Each finite test is
Schwartz--Bruhat, so the formula is also ordinary lattice Poisson summation
with congruence classes.  At \(p\notin S\), trace duality is unimodular and
\[
 \widehat\eta_{p,l}(A)
   =p^{-3v_p(l)}
       \psi_p(\Tr(Az_p^0))
       \mathbf 1_{p^{-v_p(l)}R_{F,p}}(A)
       \quad (p\mid l);
\]
for \(p\nmid l\) it is \(\mathbf 1_{R_{F,p}}\).
In particular every such transform is bounded in absolute value by one.
The finite Fourier factors in Equation~\eqref{eq:add-poisson} are supported
on the fractional ideal
\begin{equation}
 \Gamma_{l,\mathfrak h}=l^{-1}\mathfrak h^{-1},\qquad P=l^3H,
 \qquad
 \Nm(\Gamma_{l,\mathfrak h})=P^{-1}.
 \label{eq:add-frequency-lattice}
\end{equation}

The following elementary geometry of this ideal will be used both to count
frequencies and to express the local congruences in global integer
coordinates.

\begin{lemma}[Lattice bounds for a cubic ideal]
\label{lem:add-ideal-lattice}
Let \(\mathfrak a\) be a fractional ideal of a totally real cubic field of
discriminant \(D_F\), and write \(\Nm(\mathfrak a)=P^{-1}\) with \(P>0\).
Let \(B_Y=\{x\in\R^3:\|x\|\le Y\}\).
In embedding coordinates it has covolume \(\sqrt{D_F}/P\), first Euclidean
minimum \(\gg P^{-1/3}\), and, for \(Y>0\),
\begin{equation}
 \#\{A\in\mathfrak a:\|A\|\le Y\}
  \ll 1+YP^{1/3}+Y^2P^{2/3}+\frac{Y^3P}{\sqrt{D_F}} .
 \label{eq:add-lattice-count}
\end{equation}
There is a \(\Z\)-basis \(b_1,b_2,b_3\) of \(\mathfrak a\) in which every
\(A=\sum n_i b_i\) with \(\|A\|\le Y\) satisfies
\[
 |n_i|\ll YP^{1/3}\quad(1\le i\le3).
\]
All constants here are absolute.
\end{lemma}

\begin{proof}
The covolume formula follows from the discriminant.  For \(0\ne A\in
\mathfrak a\), containment \((A)\subset\mathfrak a\) implies
\(|\Nm_{F/\Q}A|\ge \Nm\mathfrak a=P^{-1}\).  The product inequality in
three real coordinates gives
\(|\Nm A|\le(\|A\|/\sqrt3)^3\), proving the first-minimum bound.

Let \(\lambda_1\le\lambda_2\le\lambda_3\) be the Euclidean successive
minima.  A lattice in fixed dimension has a basis of bounded orthogonality
defect, by Mahler's basis theorem \cite[Section~2, Theorem~1]{Mahler1946}.
For such a basis, order its lengths increasingly and write its defect as
\(\omega=\prod_i\|b_i\|/\covol(\mathfrak a)\ll1\).
Cramer's rule and Hadamard's inequality bound the norm of the \(i\)-th
coordinate functional by \(\omega/\|b_i\|\).  Since
\(\|b_i\|\ge\lambda_i\), counting its possible integer coordinates gives
\[
 \#(\mathfrak a\cap B_Y)\ll
       \prod_{i=1}^3(1+Y/\lambda_i).
\]
Minkowski's second theorem for the Euclidean ball, in the formulation of
\cite[Theorem~4.3]{Widmer2010}, gives
\(\lambda_1\lambda_2\lambda_3\asymp\covol(\mathfrak a)\).
Expanding the last product, using \(\lambda_2\ge\lambda_1\), and inserting
the first-minimum and covolume bounds proves
Equation~\eqref{eq:add-lattice-count}.  The same coordinate functional
bound with \(\|b_i\|\ge\lambda_1\) proves the final assertion.
\end{proof}

\subsection{The nonscalar contribution}

Scalar frequencies \(A\in\Q\subset F\) in
Equation~\eqref{eq:add-poisson} record only the trace.  The main work is to
show that the remaining frequencies contribute an absolute power saving.
We prove that estimate before evaluating the scalar terms.

\begin{lemma}[Nonscalar contribution]
\label{lem:add-nonscalar}
There are absolute \(\tau,\theta,\epsilon_F>0\), with
\(\epsilon_F>2\tau+3\theta\), such that for \(C_3=371\) and the data of
Proposition~\ref{prop:additive-equidistribution} with these
\(\tau,\theta\),
\begin{equation}
 \sum_{\substack{\mathfrak h\\H\le X^{C_3}}}
   \left|
    \sum_{A\in F\setminus\Q}
       \widehat\phi(A)\widehat\eta_{\mathfrak h,l}(A)
   \right|
     \ll X^{-\epsilon_F}.
 \label{eq:add-nonscalar}
\end{equation}
The constants have the same uniformity as in that proposition.
\end{lemma}

\begin{proof}
We use \(\lambda=1/3\).  Lemma~\ref{lem:ext-order-scale}, which gives
\(\Delta_F\ge X^{1/2}/2\), implies
\begin{equation}
 \Delta_F\ge X^\lambda
 \label{eq:add-order-lower}
\end{equation}
for all sufficiently large \(X\).

After bounding the real Fourier tail, we split according to the field
discriminant.  If \(D_F\) is large, comparing it with the discriminant of
the order generated by a retained nonscalar frequency after clearing its
denominators forces \(H\) large.  Otherwise the lower bound for
\(\Delta_F\) makes the index \([R_F:C_F]\) large.  For profiles with small
\(G_0\), this forces \(H\) large; inverse support and boundary decay then
bound the possibly signed original sum \(S_{\mathfrak h}\).  Poisson
expresses its full nonscalar Fourier sum as \(S_{\mathfrak h}\) minus the
separately estimated scalar series.  For profiles with large \(G_0\), the
inverse depths attached to each frequency and its chosen shifts provide
either sufficient pointwise decay or congruences; in the latter range a
further discriminant comparison supplies the required denominator size.

The positive exponents
\(d,e,e_1,e_2,w,\theta,\tau\) below will be fixed at the end of the proof;
put \(Y=X^w\) and \(s=w+\theta\).

First the frequencies outside \(\|A\|\le Y\) can be discarded.  For every
fixed integer \(r\ge0\), integration by parts in the fixed real coordinates
gives, uniformly over the box family,
\begin{equation}
 |\widehat\phi(A)|\ll_r
       X^{r\tau}(1+\|A\|)^{-r}.
 \label{eq:add-real-decay}
\end{equation}
Indeed each derivative of order \(r\) costs at most a constant times
\(X^{r\tau}\), and all supports are in one compact set.
Lemma~\ref{lem:fourier-shell-estimates} and the bounds for \(p\notin S\) give
\(|\widehat\eta_{\mathfrak h,l}(A)|\ll H^{-1}\).
The first-minimum bound also gives the rough count
\(\#(\Gamma_{l,\mathfrak h}\cap B_R)\ll(1+RP^{1/3})^3\).
For \(R\ge Y\ge1\) this is \(O(R^3P)\), since \(P=l^3H\ge1\).
Dyadic summation of Equation~\eqref{eq:add-real-decay}, for \(r>3\), yields
\begin{equation}
 \sum_{\substack{A\in F\\\|A\|>Y}}
   |\widehat\phi(A)\widehat\eta_{\mathfrak h,l}(A)|
  \ll_r X^{r\tau}l^3Y^{3-r}
  \le X^{r\tau+3\theta+w(3-r)}\,O_r(1).
 \label{eq:add-real-tail}
\end{equation}
This also proves absolute convergence in Equation~\eqref{eq:add-poisson}.
The profile count costs only \(X^{o(1)}\).  When bounding Fourier terms
directly, we may therefore restrict to the ball and use the uniform bound
\(|\widehat\phi|\ll1\) there.

Suppose first that \(D_F\ge X^d\).  A nonscalar \(A\) generates \(F\), since
the field has degree three.  Moreover \(HR_F\subset\mathfrak h\): the
integer \(H=|R_F/\mathfrak h|\) annihilates this finite group.  Thus
\(lHA\in R_F\) for \(A\in\Gamma_{l,\mathfrak h}\).  The discriminant of the
order \(\Z[lHA]\) is a positive integer multiple of \(D_F\), whereas the
product of the squared differences of its three real conjugates is
\(O((lHY)^6)\).  Existence of a nonscalar frequency in the ball therefore
implies
\[
 D_F\ll(lHY)^6,\qquad H\gg X^{d/6-s}.
\]
Apply Equation~\eqref{eq:add-lattice-count} to
\(\Gamma_{l,\mathfrak h}\) and multiply it by \(H^{-1}\).  The four resulting
terms are bounded by constants times
\begin{equation}
 X^{-d/6+s}
 +X^{-d/9+5s/3}
 +X^{-d/18+7s/3}
 +X^{-d/2+3s}.
 \label{eq:add-large-field}
\end{equation}
All have a fixed power saving once \(s<d/100\).  If the ball contains no
nonscalar frequency there is nothing to count for that profile.

We now assume \(D_F<X^d\).  When \(d\le\lambda/3\),
Equation~\eqref{eq:add-order-lower} gives
\[
 [R_F:C_F]=(\Delta_F/D_F)^{1/2}\ge X^{\lambda/3}.
\]
Profiles with \(G_0\le X^e\) are best estimated in their original form
\(S_{\mathfrak h}\).  Equation~\eqref{eq:add-profile-indices} implies
\(H\ge X^{\lambda/3-e}\).  All supporting points lie in
\(\mathfrak a_{\mathfrak h}\), whose reciprocal norm is
\(O(D_FG_0^{24})\) by Equation~\eqref{eq:add-profile-support}.
The first-minimum bound of Lemma~\ref{lem:add-ideal-lattice}, applied in the
fixed real region, gives \(O(D_FG_0^{24})\) possible points.
At each such point, Equations
\(\eqref{eq:add-profile-norm}\)--\(\eqref{eq:add-profile-inverse}\) give
\[
 \left|\prod_{p\in S}K_{p,\mathfrak h_p}(z)\right|
   \ll X^{o(1)}
        \left(\frac{D_FG_0^{24}}{H}\right)^{1/2}.
\]
Here and below \(X^{o(1)}\) denotes \(O_\eta(X^\eta)\) for every fixed
\(\eta>0\), uniformly in the arithmetic data; it includes factors
\(8^{|S|}\).  It follows that
\begin{equation}
 |S_{\mathfrak h}|
   \ll X^{-\lambda/6+(3/2)d+(73/2)e+o(1)}.
 \label{eq:add-small-index-profile}
\end{equation}
The use of the original sum has gained the boundary factor from the inverse
transform.  It remains to check that its scalar portion is small as well,
so that this bounds the nonscalar portion.

If a rational scalar has denominator \(p^k\) at \(p\), its denominator
ideal in \(R_{F,p}\) is \(p^kR_{F,p}\); its component exponents are
\(e_{\mathfrak q}k\), so its norm is \(p^{3k}\), including at ramified
primes.  Thus a scalar in this profile has its \(S\)-part denominator
\(t=H^{1/3}\), or the profile contains no scalar.  Its denominator away
from \(S\) divides \(l\), so there are
\(O(1+YlH^{1/3})\) such rationals in the ball.  Their absolute contribution
is at most
\begin{equation}
 O\bigl(H^{-1}+YlH^{-2/3}\bigr).
 \label{eq:add-small-index-scalars}
\end{equation}
By \(H\ge X^{\lambda/3-e}\) this also saves a power for the choices below.
Equation~\eqref{eq:add-poisson} expresses the full dual sum as
\(S_{\mathfrak h}\).  Subtracting its scalar frequencies in the ball and
using Equation~\eqref{eq:add-real-tail} for the scalar tail therefore bounds
the full nonscalar sum for these profiles by
Equations~\eqref{eq:add-small-index-profile}--\eqref{eq:add-small-index-scalars}
and the tail.  No positivity of an individual inverse shell is needed.

It remains to treat \(D_F<X^d\) and \(G_0>X^e\).  For each frequency in
the denominator shell of the profile choose at each \(p\in S\) one of the shifts \(s_p\) of
Lemma~\ref{lem:fourier-shift}, and write
\[
 \alpha_p=A+s_p,\qquad \beta_p=\alpha_p^{-1},\qquad
 J=\prod_{p\in S}p^{j_p},
\]
where \(j_p\) is the least nonnegative integer with
\(p^{j_p}\beta_p\in C_{F,p}\).  All sums involving \(J\) below are over
frequencies in this profile.  The pointwise local estimate now reads
\[
 |\widehat\eta_{\mathfrak h,l}(A)|
       \ll H^{-1}J^{-1/2}.
\]
For frequencies with \(J\ge X^{e_2}\), the rough lattice count immediately
gives
\begin{equation}
 \sum_{\substack{\|A\|\le Y\\J\ge X^{e_2}}}
  |\widehat\phi(A)\widehat\eta_{\mathfrak h,l}(A)|
  \ll \frac{(1+YP^{1/3})^3}{H}\,X^{-e_2/2}
  \ll X^{3s-e_2/2}.
 \label{eq:add-large-j}
\end{equation}
We used \(H,l,Y\ge1\).  This leaves \(J<X^{e_2}\), where the extra factor in
the Fourier bound may be too small.  The order discriminant then forces a
large denominator ideal, and the definition of \(j_p\) supplies congruences
among its possible frequencies.

For a nonscalar \(A\) let
\(\operatorname{disc}(A)=\Disc(1,A,A^2)\in\Q^\times\).
The equality \(F=\Q[A]\) remains
\(F\otimes\Q_p=\Q_p[A]\) after tensoring.  Shifting by a scalar and
inverting an element with nonzero components preserve the generated
algebra, so \(p^{j_p}\beta_p\) generates \(F\otimes\Q_p\).
The full-rank order \(\Z_p[p^{j_p}\beta_p]\) is contained in \(C_{F,p}\).
The identities
\[
 \operatorname{disc}(\beta_p)
   =\frac{\operatorname{disc}(A)}
          {\Nm_{F_p/\Q_p}(\alpha_p)^4},
 \qquad
 \operatorname{disc}(p^{j_p}\beta_p)
   =p^{6j_p}\operatorname{disc}(\beta_p)
\]
follow by taking the three differences of conjugates and by scaling.
Discriminant divisibility therefore gives
\[
 v_p(\operatorname{disc}(A))
  -4v_p(\Nm\alpha_p)+6j_p\ge v_p(\Delta_F).
\]
The shift lemma says \(v_p(\Nm\alpha_p)=-\log_pH_p\) for \(p>5\), and
differs from this by at most three at \(2,3,5\).
For \(p\notin S\), \(lA\) is integral and \(C_{F,p}=R_{F,p}\) is unramified,
so \(v_p(\operatorname{disc}(A))\ge-6v_p(l)\).
Combining these valuation bounds for the nonzero rational
\(\operatorname{disc}(A)\), with only an absolute factor from \(2,3,5\),
gives
\begin{equation}
 Y^6\gg |\operatorname{disc}(A)|
   \gg \Delta_F\,J^{-6}H^{-4}l^{-6}.
 \label{eq:add-frequency-discriminant}
\end{equation}
The first inequality uses \(\|A\|\le Y\).
Thus for \(J<X^{e_2}\), Equations
\(\eqref{eq:add-order-lower}\) and \(\eqref{eq:add-frequency-discriminant}\)
imply
\begin{equation}
 H\gg X^{\lambda/4-(3/2)(s+e_2)}.
 \label{eq:add-denominator-lower}
\end{equation}
In particular, once \(s+e_2<\lambda/12\) with a margin, every nonscalar
frequency under consideration has \(H\ge X^{\lambda/8}\) for large \(X\).

For the remaining count, group frequencies according to their \(j_p\)'s and
the chosen \(s_p\)'s.  The number of exponent lists with \(J<X^{e_2}\) is
\(X^{o(1)}\), by the same Rankin argument as for profiles; the bounded lists
of shifts cost \(4^{|S|}=X^{o(1)}\).  Groups containing no nonscalar
frequency contribute zero and are omitted.  Fix one of the other groups.
Apply the basis
part of Lemma~\ref{lem:add-ideal-lattice} to \(\Gamma_{l,\mathfrak h}\).
All frequencies in the ball have integer coordinates in a cube
\([-L,L]^3\), where we choose \(L\) comparable to the upper bound
\(YP^{1/3}\).  This choice, together with
Equation~\eqref{eq:add-denominator-lower}, gives
\begin{equation}
 L\asymp YP^{1/3},\qquad
 L\gg X^{\lambda/24},\qquad
 L^3/H\ll X^{3s}.
 \label{eq:add-coordinate-cube}
\end{equation}

At \(p\in S\) choose a generator \(i_p\) of the maximal-order ideal
\(\mathfrak h_p\), and set \(x=i_p(A+s_p)\).  Localizing the global basis
of \(\Gamma_{l,\mathfrak h}\) at \(p\), then multiplying it by \(i_p\), gives
a \(\Z_p\)-basis of \(R_{F,p}\), because \(p\nmid l\).  Accordingly the
map from the three integer coordinates to \(x\) is an affine
\(\operatorname{GL}_3(\Z_p)\) change of coordinates.  This statement holds
also for a product of ramified local fields.

Smith normal form for \(\mathfrak h_p\cap C_{F,p}\) inside \(\mathfrak h_p\)
gives a primitive linear form \(\ell_p:R_{F,p}\to\Z_p\), meaning that its
image is \(\Z_p\), and an integer
\(g_p'\ge g_p/3\) such that
\[
 \ell_p\bigl(i_p^{-1}(\mathfrak h_p\cap C_{F,p})\bigr)
       \subset p^{g_p'}\Z_p.
\]
Indeed the three elementary-divisor exponents sum to \(g_p\), and we use
the coordinate with largest exponent.  If \(p>5\), the shift lemma makes
\(x\) a unit and \(\beta_p=i_px^{-1}\in\mathfrak h_p\).  The condition
\(p^{j_p}\beta_p\in C_{F,p}\) implies
\begin{equation}
 \ell_p(x^{-1})\equiv0\pmod {p^k}
       \quad\text{for every }1\le k\le(g_p'-j_p)_+.
 \label{eq:add-inverse-congruence}
\end{equation}
Inversion permutes the units in the finite algebra
\(R_{F,p}/p^kR_{F,p}\).  A primitive linear form has \(p^{2k}\) zeros among
the \(p^{3k}\) additive residues.  Thus the unit condition and
Equation~\eqref{eq:add-inverse-congruence} retain at most the fraction
\(p^{-k}\) of all residues modulo \(p^k\).  Nilpotents in the finite algebra
do not affect this argument.

For \(p=2,3,5\), the shift lemma gives \(x\in R_{F,p}\) and
\(x^{-1}\in p^{-1}R_{F,p}\).  Put \(y=p x^{-1}\).  Then
\(p^{j_p}i_py=p(p^{j_p}\beta_p)\) still lies in
\(\mathfrak h_p\cap C_{F,p}\), and the congruence in
Equation~\eqref{eq:add-inverse-congruence} holds for \(y\) in place of
\(x^{-1}\).  On the permitted locus the inverse identities
\[
 y(x)-y(x')=p x^{-1}(x'-x)(x')^{-1},\qquad
 x-x'=p^{-1}x\bigl(y(x')-y(x)\bigr)x'
\]
show respectively that \(x\bmod p^{k+1}\) determines \(y\bmod p^k\), and
that a fixed residue of \(y\bmod p^k\) confines \(x\) to one residue modulo
\(p^{k-1}\).  The permitted locus has \(v_{\mathfrak q}(x)=0\) on the pole
components and \(0\le v_{\mathfrak q}(x)\le e_{\mathfrak q}\) on the other
components.  These conditions are determined modulo \(p^2\), since a
change in \(p^2R_{F,p}\) has component valuation at least \(2e_{\mathfrak q}\).
At modulus \(p^{k+1}\), each of the at most \(p^{2k}\) residues in
the linear kernel therefore allows at most \(p^6\) residues of \(x\).
The retained fraction is at most \(p^{3-k}\), hence \(O(p^{-k})\) with an
absolute constant at these three primes.  This proves the required
congruence fraction with only one extra power of a small prime in its
modulus.

Write \(a_p=(g_p'-j_p)_+\).  The available congruence powers satisfy
\begin{equation}
 \prod_{p\in S}p^{a_p}
   \ge G_0^{1/3}/J
   > X^{e/3-e_2}.
 \label{eq:add-available-congruences}
\end{equation}
We will choose \(e/3-e_2>2e_1\).  If every prime with \(a_p>0\) is at most
\(X^{e_1}\), multiply available factors \(p\), one at a time, until their
product \(M_1\) first reaches \(X^{e_1/2}\).  Equation
\(\eqref{eq:add-available-congruences}\) ensures this is possible, and
\[
 X^{e_1/2}\le M_1\le X^{3e_1/2}<X^{2e_1}.
\]
Impose just the corresponding congruences.  The Chinese remainder theorem
combines their retained fractions into \(O(M_1^{-1})\).  Their joint modulus
is \(M_1\) times at most \(2\cdot3\cdot5\), because only the small-prime
tests need one extra power.  The affine integral coordinate changes above
preserve these fractions in the global coordinates.  If
\(2e_1<\lambda/24\), this joint modulus is \(o(L)\) by
Equation~\eqref{eq:add-coordinate-cube}.  Counting each admissible residue
class in the cube then gives
\begin{equation}
 O(L^3/M_1)=O(L^3X^{-e_1/2})
 \label{eq:add-composite-congruence-count}
\end{equation}
possible frequencies.

Otherwise some available prime satisfies \(p>X^{e_1}\); for large \(X\)
it exceeds \(5\).  Use its congruence modulo \(p\).  Set
\(B=R_{F,p}/pR_{F,p}\), a three-dimensional commutative \(\mathbf F_p\)-algebra,
possibly with nilpotents, and reduce \(\ell_p\) modulo \(p\).  Choose an
\(\mathbf F_p\)-basis of \(B\) and let \(x\) be the universal element of
\(B\otimes_{\mathbf F_p}\mathbf F_p[X_1,X_2,X_3]\) in this basis.
Let \(M_x\) be its multiplication matrix and put
\(\mathcal N_B(x)=\det M_x\).  This is a nonzero polynomial of degree three,
since \(\mathcal N_B(1)=1\).  On the unit open set \(\mathcal N_B\ne0\),
\[
 \ell_p(x^{-1})=\frac{P_\ell(x)}{\mathcal N_B(x)},\qquad
 P_\ell(x)=\ell_p\bigl(\operatorname{adj}(M_x)1_B\bigr),
\]
where \(P_\ell\) has degree at most two.  This numerator is nonzero.
Indeed the affine space parametrizing elements of the vector space \(B\)
has coordinate ring \(\mathbf F_p[X_1,X_2,X_3]\), which is a domain even
when multiplication in \(B\) has nilpotents.  The unit open set is dense,
and inversion is an automorphism of that open set.  The nonzero linear
polynomial \(\ell_p\) stays nonzero on the open set and after composition
with inversion.  The displayed formula, or injectivity of localization at
\(\mathcal N_B\), then forces \(P_\ell\ne0\).
The affine \(\operatorname{GL}_3(\mathbf F_p)\) change from global lattice
coordinates preserves nonvanishing.

Let \(N_L=2\lfloor L\rfloor+1\asymp L\) be the number of possible integers in
one coordinate of the cube.  A nonzero polynomial of bounded degree has
\(O(N_L^2)\) zeros on a product of three residue intervals of length \(N_L<p\);
this follows by induction on the number of variables, viewing it as a
polynomial in the last variable and bounding the slices on which all its
coefficients vanish.  If \(N_L\ge p\), partition the cube into boxes of side
at most \(p\) and use the same bound there.  The number of global integer
coordinates satisfying \(P_\ell=0\bmod p\) is consequently
\begin{equation}
 O\bigl(L^3/p+L^2\bigr)
   =O\bigl(L^3(X^{-e_1}+X^{-\lambda/24})\bigr).
 \label{eq:add-prime-congruence-count}
\end{equation}
This bounds the subset for which \(x\) is a unit as well.
Multiplying Equations
\(\eqref{eq:add-composite-congruence-count}\)--\(\eqref{eq:add-prime-congruence-count}\)
by \(H^{-1}\), and using Equation~\eqref{eq:add-coordinate-cube}, gives
respectively the bounds
\[
 O(X^{3s-e_1/2}),\qquad
 O(X^{3s-e_1}+X^{3s-\lambda/24})
\]
for the fixed group of \(j_p\)'s and shifts.  Their \(X^{o(1)}\) number is
harmless once these exponents are negative with a margin.  This completes
the count of nonscalar frequencies in the ball.

We finish by making the parameter choices and all margins explicit.  Put
\(\lambda_0=\min(\lambda,1)\) and take
\begin{equation}
 \begin{aligned}
 d=e&=\frac{\lambda_0}{1000},&
 e_1&=\frac e{100},&
 e_2&=\frac{e_1}{100},\\
 w&=\frac{e_2}{100},&
 \theta=\tau&=\frac w{100},&
 \epsilon_F&=\frac{e_2}{4}.
 \end{aligned}
 \label{eq:add-parameters}
\end{equation}
Then \(s=101e_2/10000\),
\[
 s<d/100,\quad d<\lambda/3,\quad
 e/3-e_2>2e_1,\quad 2e_1<\lambda/24,\quad
 s+e_2<\lambda/12.
\]
The saving in Equation~\eqref{eq:add-small-index-profile} is at least
\[
 \lambda/6-(3/2)d-(73/2)e
    \ge (1/6-0.038)\lambda_0,
\]
and Equation~\eqref{eq:add-small-index-scalars} has a still larger margin
than \(\epsilon_F\), since \(e\le\lambda/1000\).
All four exponents in Equation~\eqref{eq:add-large-field} also have margins
at least \(29d/900>\epsilon_F\), by \(s<d/100\).
Among the remaining ball estimates the smallest
margin is in Equation~\eqref{eq:add-large-j}, namely
\[
 e_2/2-3s=\frac{4697}{10000}e_2>\epsilon_F;
\]
the congruence margins \(e_1/2-3s\), \(e_1-3s\), and
\(\lambda/24-3s\) are larger.  With \(r=100\),
the exponent in Equation~\eqref{eq:add-real-tail} is
\[
 r\tau+3\theta+w(3-r)=-\frac{9597}{10000}e_2.
\]
All \(X^{o(1)}\) costs from profiles, shifts, and exponent lists can thus be
absorbed while retaining \(\epsilon_F=e_2/4\).  This proves
Equation~\eqref{eq:add-nonscalar}.  Finally
\(2\tau+3\theta=e_2/2000<\epsilon_F\), as asserted.
\end{proof}

\subsection{The scalar term and completion}

\begin{proof}[Proof of Proposition~\ref{prop:additive-equidistribution}]
Use the constants in Equation~\eqref{eq:add-parameters}.
Equations~\eqref{eq:add-original-profiles} and \(\eqref{eq:add-poisson}\),
together with Lemma~\ref{lem:add-nonscalar}, leave the scalar frequencies
with absolute error \(O(X^{-\epsilon_F})\).  We first remove the profile
cutoff from those frequencies.

For \(a\in\Q\), let \(t=\prod_{p\in S}p^{\max(0,-v_p(a))}\) be its
\(S\)-part denominator, with \(t=1\) for \(a=0\).  As noted above, its
unique profile has \(H=t^3\).
Lemma~\ref{lem:fourier-scalar-density} gives the exact product
\(\prod_{p\in S}\gamma_p(a)=t^{-3}\), including at \(2,3,5\).
The other denominator divides \(l\), so all such scalars with a given \(t\)
lie in \((tl)^{-1}\Z\).  The change of real variables
\((z_i,z_j,z_k)\mapsto(z_i,z_j,\Tr z)\) has absolute Jacobian one.
Writing \(\mathbf 1=(1,1,1)\), it gives the exact identity
\[
 \widehat\phi(a\mathbf 1)=A_\phi\,\widehat\chi(a).
\]
For any fixed \(r>1\), rapid decay of the fixed function \(\chi\) implies
\(\sum_{n\in\Z}|\widehat\chi(n/(tl))|\ll_r tl\).
The scalars omitted by \(H\le X^{C_3}\) therefore contribute at most
\begin{equation}
 \sum_{t>X^{C_3/3}}t^{-3}
       \sum_{n\in\Z}|\widehat\phi((n/(tl))\mathbf 1)|
    \ll A_\phi\,l\,X^{-C_3/3}.
 \label{eq:add-scalar-tail}
\end{equation}
We have enlarged the set of \(t\)'s to all positive integers, which is an
upper bound.  With \(C_3=371\), this is \(O(X^{-\epsilon_F})\), since
\(A_\phi\ll1\), \(l\le X^\theta\), and \(371/3-\theta>\epsilon_F\).
The same calculation without the cutoff proves
absolute convergence of the unrestricted scalar series.

That series is one-dimensional Poisson summation for the trace marginals.
At infinity the marginal is \(A_\phi\chi(t)\).  At \(p\in S\) it is
\(d_p(t)\mathbf 1_{\Z_p}(t)\).  At \(p\mid l\), the onto integral trace map
and two kernel coordinates show that the indicator
\(z_p^0+p^{v_p(l)}R_{F,p}\) has trace density
\[
 p^{-2v_p(l)}\mathbf 1_{m+p^{v_p(l)}\Z_p}(t)
\]
with respect to Haar measure of mass one on \(\Z_p\).  At every other
finite prime the marginal is \(\mathbf 1_{\Z_p}\).
Their Fourier transforms at a rational \(a\) are exactly the finite factors
in Equation~\eqref{eq:add-poisson} evaluated at the scalar \(a\), because
\(\Tr(az)=a\Tr z\).

For completeness, this one-dimensional application of Poisson does not
require a smoothness assumption on \(d_p\).  Its Fourier coefficients on
\(\Q_p/\Z_p\) of denominator \(p^k\) are \(p^{-3k}\), and there are
\(p^k-p^{k-1}\) such characters.  Their absolute sum is finite.
Truncating each expansion gives a compactly supported locally constant
density on \(\Z_p\), to which Poisson applies.  The densities converge
uniformly by Lemma~\ref{lem:fourier-scalar-density}, and the preceding
absolute scalar bound permits passage to the limit on the Fourier side.
On the original side, the finite marginals force the rational trace to be
an integer.  Only \(m\) meets the real cutoff.  Hence the unrestricted
scalar contribution is precisely
\[
 A_\phi\,l^{-2}\prod_{p\in S}d_p(m).
\]

Lemma~\ref{lem:fourier-scalar-density} bounds the product of \(d_p(m)\)
above and below by positive absolute constants.  Also
\(A_\phi\gg X^{-2\tau}\) by Equation~\eqref{eq:add-box}, and
\(l^{-2}\ge X^{-2\theta}\).  Dividing the combined absolute error
\(O(X^{-\epsilon_F})\) by the main term therefore costs at most
\(O(X^{2\tau+2\theta})\).  Since
\(\epsilon_F>2\tau+3\theta\), the relative error is \(O(X^{-\theta})\).
Every estimate used fixed real derivative and support bounds and constants
independent of the fields, orders, residues, and box locations and scales.
This proves Equation~\eqref{eq:add-equidistribution} with its stated
uniformity.
\end{proof}
